\begin{table}[htbp]
\centering
\caption{Distribuição dos pares candidatos por faixa do escore do \textit{linkage} probabilístico}
\label{tab:score-bands}
\begin{tabular}{lrrr}
\toprule
Faixa do escore prob. & Total & Pares & \% Pares \\
\midrule
    0-3 & 5,336 & 0 & 0.00\% \\
    3-5 & 34,645 & 3 & 0.01\% \\
    5-6 & 14,263 & 17 & 0.12\% \\
    6-7 & 6,175 & 57 & 0.92\% \\
    7-8 & 1,080 & 43 & 3.98\% \\
    8-9 & 102 & 34 & 33.33\% \\
    9-10 & 44 & 43 & 97.73\% \\
    10+ & 51 & 50 & 98.04\% \\
\bottomrule
\end{tabular}

\smallskip
{\footnotesize\textit{Nota}: o escore do \textit{linkage} probabilístico corresponde à soma das razões de log-verossimilhança calculadas pelo OpenRecLink para três campos de comparação, a saber, nome, nome da mãe e data de nascimento (cf.\ Seção~\ref{sec:fontes-dados}). O piso ($-12{,}85$) indica discordância em todos os campos. A análise exploratória dos 116 valores distintos permitiu estimar a contribuição aproximada de cada campo quando concordante: data de nascimento ($\approx 2{,}3$ pontos), nome ($\approx 3{,}2$) e nome da mãe ($\approx 4{,}9$). Assim, escores nas faixas intermediárias (5--7) tipicamente refletem concordância parcial em dois campos, ao passo que valores acima de 9 indicam concordância simultânea na totalidade dos campos.}
\end{table}
